\documentclass[leqno]{article}
% Shared QBT/Soln pack preamble for resources.danieljsmith.org.
\usepackage{graphicx}
\usepackage{dsfont}
\usepackage{mathtools}
\usepackage{amssymb}
\usepackage{amsmath}
\usepackage{mleftright}
\usepackage{amsthm}
\usepackage{float}
\usepackage{ragged2e}
\usepackage{tensor}
\usepackage{fancyhdr}
\usepackage{empheq}
\usepackage[most]{tcolorbox}
\usepackage{enumitem}
\usepackage{dirtytalk}
\usepackage{marginnote}
\usepackage{microtype}
\usepackage[
  a4paper,
  left=0.8in,
  right=1in,
  top=1in,
  bottom=0.5in,
  headheight=14pt,
  headsep=0.4in
]{geometry}
\setlength{\parindent}{0cm}

% TikZ / PGFPlots
\usepackage{pgfplots}
\pgfplotsset{compat=1.18}
\usepgfplotslibrary{fillbetween, polar}
\usetikzlibrary{calc, positioning}
\usetikzlibrary{angles, quotes}
\usetikzlibrary{arrows.meta}
\usetikzlibrary{decorations.pathreplacing, decorations.markings}
\usetikzlibrary{patterns}
\usetikzlibrary{intersections}

% Links and contents. The contents follow the TMUA papers: compact entries,
% no dot leaders, and each question's mark total in grey after its label.
\usepackage{tocloft}
\tocloftpagestyle{fancy}
\setlength{\cftbeforesubsecskip}{2pt}
\setlength{\cftsubsecindent}{1.5em}
\setlength{\cftsubsecnumwidth}{0pt}
\renewcommand{\cftsubsecfont}{\small\color{black}}
\renewcommand{\cftsubsecpagefont}{\small\color{black}}
\renewcommand{\cftsubsecleader}{\hfill}
\usepackage[colorlinks=true,linkcolor=black,citecolor=magenta,urlcolor=black]{hyperref}%
\urlstyle{same}

% Pack macros
% \marks also adds its value to the current question's total in the .aux
% file, so the contents show each total from the second pass. The write
% happens at shipout, so a box typeset more than once counts once. The
% pack's mark total for the contents-page title is summed the same way.
\makeatletter
\newcount\djs@totalmarks
\newcommand{\djsAddMarks}[2]{%
  \global\advance\djs@totalmarks#2\relax
  \@ifundefined{djs@marks@#1}%
    {\expandafter\xdef\csname djs@marks@#1\endcsname{#2}}%
    {\expandafter\xdef\csname djs@marks@#1\endcsname
      {\the\numexpr\csname djs@marks@#1\endcsname+#2\relax}}}
\RenewDocumentCommand{\marks}{o m}{%
  \protected@write\@auxout{}{\string\djsAddMarks{\number\value{enumi}}{#2}}%
  \IfValueTF{#1}%
    {\marginnote{\raggedright\textbf{[#2]}}[#1]}%
    {\marginnote{\raggedright\textbf{[#2]}}}%
}
\makeatother

\newcommand{\vspacerule}[2]{%
  \par\vspace*{#1}%
  \noindent\hrulefill\par
  \vspace*{#2}%
}

\definecolor{scarlet}{RGB}{205, 40, 75}
% Site colours for the contents-page title: the accent, the rule colour, a
% stronger rule colour for the path separator, and the muted text colour.
\definecolor{djsBurgundy}{RGB}{118, 59, 67}
\definecolor{djsHairline}{RGB}{216, 213, 205}
\definecolor{djsSlash}{RGB}{170, 166, 158}
\definecolor{djsMuted}{RGB}{96, 101, 106}

% Worked-solution box, adapted from the TMUA worked-solutions box. A white breakable box with a left scarlet rule and no surrounding
% frame or tint. A part that breaks fills the rest of its page, so the rule
% runs to the bottom margin; the unbroken or last part ends in a short
% horizontal foot rule so a solution visibly stops. Unlike the TMUA box there
% is no answer label. The box does not grow to the left, so the rule sits at
% the text edge; growing to the right keeps the usable line width close to the
% old tcolorbox Solution box.
% Spacing inside the box: paragraphs are 6pt apart, so source blank lines show
% (solutions put one after each display and none between a lead-in and its
% display). Every display gets 6pt above and 10pt below, whatever the length
% of its lead-in line; align rows are spread by an 8pt \jot; solution parts
% are 14pt apart and sub-parts 8pt; \left and \right space like ordinary
% brackets (mleftright); no single line is stranded at a page break; and
% neither the Solution label nor the line before a list ends a page. No
% skip can shrink, so full pages are not squeezed.
\newtcolorbox{djsSolution}{%
  enhanced, breakable, frame hidden, sharp corners, colback=white,
  height fixed for=first and middle, height fill,
  extras unbroken and last={natural height},
  borderline west={2.2pt}{0pt}{scarlet},
  left=11pt, right=0pt, top=5pt, bottom=13pt,
  before skip=13pt, after skip=4pt,
  grow to right by=10mm,
  overlay unbroken and last={%
    \fill[scarlet] (frame.south west)
      rectangle ([xshift=3.5cm, yshift=2.2pt]frame.south west);},
  before upper={{\color{scarlet}\bfseries Solution}\par\nobreak\vspace{4pt}%
    \setlength{\parskip}{6pt plus 2pt}%
    \setlength{\abovedisplayskip}{6pt plus 2pt}%
    \setlength{\abovedisplayshortskip}{6pt plus 2pt}%
    \setlength{\belowdisplayskip}{10pt plus 2pt}%
    \setlength{\belowdisplayshortskip}{10pt plus 2pt}%
    \setlength{\jot}{8pt}%
    \setlist{parsep=6pt plus 2pt, topsep=6pt plus 2pt, beginpenalty=10000}%
    \setlist[enumerate,1]{itemsep=14pt plus 2pt}%
    \setlist[enumerate,2]{itemsep=8pt plus 2pt}%
    \mleftright
    \clubpenalty=10000 \widowpenalty=10000 \displaywidowpenalty=10000\relax}}

% Optional smaller figure inside a solution box. \scalebox shrinks node text
% with the drawing. Not applied to existing figures.
\newcommand{\SolnFigureScale}{0.85}
\newsavebox{\djsSolnFigureBox}
\newenvironment{djsSolutionFigure}
  {\par\vspace{4pt}\begin{center}\begin{lrbox}{\djsSolnFigureBox}}
  {\end{lrbox}\scalebox{\SolnFigureScale}{\usebox{\djsSolnFigureBox}}%
   \end{center}\vspace{2pt}\par}

\newcommand{\cxl}[2]{%
  \tikz[baseline=(X.base)]{%
    \node[inner sep=0pt, outer sep=0pt, text=#1] (X) {$#2$};
    \draw[#1, line width=0.4pt] (X.south west)--(X.north east);
  }%
}

% Page-style hook run at the contents page and each question's opening page.
% A no-op for QBT packs; \djsSolnHeader points it at the fancy header.
\newcommand{\djsQuestionPageStyle}{}

% A contents entry: "Question N" in a box wide enough for two digits, so the
% mark totals line up, then the total in grey. Bookmarks keep "Question N".
\makeatletter
\newlength{\djsTocLabelWidth}
\newcommand{\djsTocMarks}[1]{%
  \@ifundefined{djs@marks@#1}{}{%
    {\footnotesize\itshape\color{black!45}%
     \csname djs@marks@#1\endcsname~mark%
     \ifnum\csname djs@marks@#1\endcsname=1 \else s\fi}}}
\newcommand{\djsTocEntry}[1]{%
  \texorpdfstring
    {\settowidth{\djsTocLabelWidth}{Question 00}%
     \makebox[\dimexpr\djsTocLabelWidth+0.9em\relax][l]{Question #1}%
     \djsTocMarks{#1}}%
    {Question #1}}
\makeatother

% The question count for the contents-page title is written to the .aux at
% the end of the document, so it also appears from the second pass.
\newcounter{djsQuestions}
\newcommand{\questionitem}{%
  \djsQuestionPageStyle
  \item\phantomsection
  \stepcounter{djsQuestions}%
  \addcontentsline{toc}{subsection}{\protect\djsTocEntry{\number\value{enumi}}}%
}
\makeatletter
\newcommand{\djsSetQuestionCount}[1]{\gdef\djs@questioncount{#1}}
\AtEndDocument{\immediate\write\@auxout
  {\string\djsSetQuestionCount{\number\value{djsQuestions}}}}
\makeatother

\renewcommand{\contentsname}{Questions}

% \djsPackHeader[<link content tag>]{<left header>}; the tag says which
% document a click on the header link came from.
\newcommand{\djsPackHeader}[2][questions]{%
  \pagestyle{fancy}%
  \fancyhead{}%
  \fancyfoot{}%
  \renewcommand{\headrulewidth}{0.9pt}%
  \lhead{#2}%
  \chead{}%
  \rhead{\href{https://danieljsmith.org/?utm_source=pdf&utm_campaign=further-maths&utm_content=#1}{danieljsmith.org}}%
}

% The header's topic is also the contents-page title, under which a label
% says which booklet this is.
\newcommand{\djsPackTitle}{}
\newcommand{\djsBookletLabel}{}
\newcommand{\djsQbtHeader}[1]{%
  \djsPackHeader{#1}%
  \renewcommand{\djsPackTitle}{#1}%
  \renewcommand{\djsBookletLabel}{Questions by topic}%
}

% The colophon at the foot of page one, independent of the page style: the
% contents-page title's hairline led by a short scarlet rule, then the date on
% the left and a link to the site on the right, all inside the bottom margin.
% Its sizes are fixed, so it matches the TMUA papers' colophon. The document
% supplies the date; recompiling never changes it.
\makeatletter
\newcommand{\djsLastUpdated}[1]{%
  \AtBeginDocument{\AddToHookNext{shipout/foreground}{\djs@colophon{#1}}}}
\newcommand{\djs@colophon}[1]{%
  \put(0.8in,-\dimexpr\paperheight-0.47in\relax){%
    \makebox[0pt][l]{\color{djsHairline}%
      \rule[0.9pt]{\dimexpr\paperwidth-1.8in\relax}{0.6pt}}%
    \makebox[0pt][l]{\color{scarlet}\rule{3.5cm}{2.2pt}}}%
  \put(0.8in,-\dimexpr\paperheight-0.32in\relax){\makebox[0pt][l]{%
    \fontsize{9}{11}\selectfont Last updated #1}}%
  \put(\dimexpr\paperwidth-1in\relax,-\dimexpr\paperheight-0.32in\relax){%
    \makebox[0pt][r]{\fontsize{7}{8.5}\selectfont\color{djsMuted}%
      Latest version at \href{https://resources.danieljsmith.org/}%
        {\color{black}resources.danieljsmith.org}}}}
\makeatother

% Solutions packs show the header only on the contents page and each
% question's opening page, so it never appears part-way through a solution.
% \thispagestyle marks the next page shipped out, which is the question's
% page because every solution question starts after a \newpage.
\newcommand{\djsSolnHeader}[1]{%
  \djsPackHeader[solutions]{\textit{Solutions} - #1}%
  \pagestyle{empty}%
  \renewcommand{\djsQuestionPageStyle}{\thispagestyle{fancy}}%
  \renewcommand{\djsPackTitle}{#1}%
  \renewcommand{\djsBookletLabel}{Worked solutions}%
}

% \djsSitePath{<strand>}{<topic>}, called once after the header, names where
% the pack sits on the site; the contents-page title prints it above the pack
% title, e.g. CORE PURE / FURTHER CALCULUS. Without it the line is left out.
\makeatletter
\newif\ifdjs@sitepath
\newcommand{\djs@strand}{}
\newcommand{\djs@section}{}
\newcommand{\djsSitePath}[2]{%
  \djs@sitepathtrue
  \renewcommand{\djs@strand}{#1}%
  \renewcommand{\djs@section}{#2}%
}
\makeatother

% Further Mechanics packs call \djsMechanicsPack in their preamble, so that
% \djsFrontMatter prints the default value of g above the contents. The box
% keeps its own space above it (before skip is the whole gap below the title,
% not an addition to it). Without the box, \djsGravityNote is just extra space
% before the contents heading, bringing the gap above it close to the gap
% between the box and the heading on Further Mechanics pages.
\newcommand{\djsGravityNote}{\vspace{14pt}}
\newcommand{\djsMechanicsPack}{%
  \renewcommand{\djsGravityNote}{%
    \begin{tcolorbox}[sharp corners, colback=white, colframe=black,
      boxrule=1pt, left=8pt, right=8pt, top=6pt, bottom=6pt,
      before skip=20pt]
    {\bfseries Instructions}\par\vspace{4pt}
    Unless a question states otherwise, take $g = 9.8\,\mathrm{m\,s^{-2}}$
    wherever a numerical value of $g$ is required. Solutions also give final
    answers for $g = 10\,\mathrm{m\,s^{-2}}$ where these differ.
    \end{tcolorbox}%
    \vspace{16pt}}%
}

% Contents-page title: the site path in spaced burgundy capitals, the pack
% title, a hairline led by a short scarlet rule (the solution box's foot rule),
% then the booklet label with the question and mark counts. The counts come
% from the .aux, so they appear from the second pass.
\makeatletter
\newsavebox{\djs@eyebrowbox}
\newcommand{\djs@caps}[1]{\textls[180]{\MakeUppercase{#1}}}
\newcommand{\djs@eyebrow}{%
  \color{djsBurgundy}\footnotesize\bfseries
  \djs@caps{\djs@strand}%
  {\color{djsSlash}\hspace{0.55em}/\hspace{0.55em}}%
  \djs@caps{\djs@section}}
\newcommand{\djs@dot}{\hspace{0.45em}\textperiodcentered\hspace{0.45em}}
\newcommand{\djs@counts}{%
  \@ifundefined{djs@questioncount}{}{%
    \djs@questioncount~question\ifnum\djs@questioncount=1 \else s\fi
    \ifnum\djs@totalmarks>0
      \djs@dot\the\djs@totalmarks~mark\ifnum\djs@totalmarks=1 \else s\fi
    \fi}}
\newcommand{\djsContentsTitle}{%
  \ifdjs@sitepath
    \sbox\djs@eyebrowbox{\djs@eyebrow}%
    \ifdim\wd\djs@eyebrowbox>\linewidth
      \@latex@warning@no@line{The \string\djsSitePath\space line is wider
        than the text, so it wraps. Shorten its topic}%
    \fi
    {\raggedright\djs@eyebrow\par}%
    \vspace{7pt}%
  \fi
  {\raggedright\fontsize{24.88}{28}\selectfont\djsPackTitle\par}%
  \vspace{9pt}%
  \noindent\makebox[0pt][l]{\color{djsHairline}\rule[0.9pt]{\linewidth}{0.6pt}}%
  {\color{scarlet}\rule{3.5cm}{2.2pt}}\par
  \vspace{5pt}%
  {\small\color{djsMuted}\djsBookletLabel\hfill\djs@counts\par}%
  \vspace{8pt}%
}
\makeatother

\newcommand{\djsFrontMatter}{%
  \djsQuestionPageStyle
  \vspace*{6pt}%
  \djsContentsTitle
  \djsGravityNote
  \tableofcontents
  \newpage
}


\djsQbtHeader{Motion Under a Variable Force}
\djsSitePath{Further Mechanics}{Kinematics \& Dynamics}
\djsMechanicsPack

\djsLastUpdated{11 October 2026}
\begin{document}
\djsFrontMatter
\begin{enumerate}[
  label=\textbf{\arabic*.},
  leftmargin=*,
  topsep=0pt,
  partopsep=0pt,
  parsep=0pt
]

% ---- Question 1 ----
\questionitem
% [Source: _QBT__Motion_Under_a_Variable_Force, Q1]

A particle $P$ of mass $4\,\mathrm{kg}$ moves along a straight horizontal line. At time $t\,\mathrm{s}$, its displacement in the positive direction from a fixed point $O$ is $x\,\mathrm{m}$ and its velocity is $v\,\mathrm{m\,s^{-1}}$. \\

The only horizontal force acting on $P$ has magnitude
\[
\dfrac{36}{(x+2)^3}\,\mathrm{N}
\] \\
and acts in the direction of decreasing $x$. Initially, $P$ is at $O$ and is moving in the positive direction with speed $\dfrac{3}{2}\,\mathrm{m\,s^{-1}}$. \\

\begin{enumerate}[label=\textbf{(\alph*)}]
  \item Find an expression for $v$ in terms of $x$. \marks{3} \\
  \item Hence find an expression for $x$ in terms of $t$. \marks{2} \\
\end{enumerate}

\newpage

% ---- Question 2 ----
\questionitem
% [Source: _QBT__Motion_Under_a_Variable_Force, Q2]

A lorry of mass $2000\,\mathrm{kg}$ travels along a straight horizontal road. Its engine delivers constant power of $80\,\mathrm{kW}$. When the lorry is moving at speed $v\,\mathrm{m\,s^{-1}}$, the resistance to its motion has magnitude $200v\,\mathrm{N}$. \\

At time $t$ seconds, the speed of the lorry is $v\,\mathrm{m\,s^{-1}}$. At $t=0$, the speed of the lorry is $8\,\mathrm{m\,s^{-1}}$. \\

\begin{enumerate}[label=\textbf{(\alph*)}]
  \item Show that
  \[
  10v\dfrac{\mathrm{d}v}{\mathrm{d}t}=400-v^2
  \]
  \marks[-30pt]{3}
  \item Verify that, while $8\leq v<20$,
  \[
  t=5\ln\left(\dfrac{336}{400-v^2}\right)
  \]
  \marks[-30pt]{5}
\end{enumerate}

\newpage

% ---- Question 3 ----
\questionitem
% [Source: _QBT__Motion_Under_a_Variable_Force, Q3]

A particle $P$ of mass $0.5\,\mathrm{kg}$ moves along the positive $x$-axis under the action of a single force. At time $t$ seconds, the particle is $x$ metres from the origin $O$, where $x>0$, and has speed $v\,\mathrm{m\,s^{-1}}$. The force acts in the direction of increasing $x$ and has magnitude
\[
\frac{k}{x^3}\text{ newtons}
\] \\
where $k$ is a positive constant. \\

When $x=2$, the speed of $P$ is $1\,\mathrm{m\,s^{-1}}$, and when $x=4$, its speed is $2\,\mathrm{m\,s^{-1}}$. \\

\begin{enumerate}[label=\textbf{(\alph*)}]
  \item Show that
  \[
  v^2=5-\frac{16}{x^2}
  \]
  \marks[-20pt]{6}
\end{enumerate}
The time taken for the speed of $P$ to increase from $1\,\mathrm{m\,s^{-1}}$ to $2\,\mathrm{m\,s^{-1}}$ is $T$ seconds. \\
\begin{enumerate}[label=\textbf{(\alph*)}, resume]
  \item Use algebraic integration to show that
  \[
  T=\frac{6}{5}
  \]
  \marks[-20pt]{4}
\end{enumerate}

\newpage

% ---- Question 4 ----
\questionitem
% [Source: _QBT__Motion_Under_a_Variable_Force, Q4]

A particle $P$ of mass $3\,\mathrm{kg}$ moves along the positive $x$-axis. At time $t=0$, it passes through the origin with speed $4\,\mathrm{m\,s^{-1}}$. \\

While $P$ is moving, the only force acting on it is a resistive force of magnitude
\[
\frac{6}{v}\,\mathrm{N}
\] \\
where $v\,\mathrm{m\,s^{-1}}$ is its speed. The force acts in the negative $x$-direction. \\

\begin{enumerate}[label=\textbf{(\alph*)}]
  \item Find an expression for $v$ in terms of the displacement $x$ of $P$ from the origin. \marks{5} \\
  \item Find the distance travelled by $P$ while its speed decreases from $4\,\mathrm{m\,s^{-1}}$ to $2\,\mathrm{m\,s^{-1}}$. \marks{2} \\
\end{enumerate}
A particle $Q$, also of mass $3\,\mathrm{kg}$, starts from the origin with speed $4\,\mathrm{m\,s^{-1}}$. The only force acting on $Q$ is a constant resistive force of magnitude $\dfrac{3}{2}\,\mathrm{N}$. \\
\begin{enumerate}[label=\textbf{(\alph*)}, resume]
  \item By comparing the resistive forces acting on $P$ and $Q$, explain why $P$ must come to rest at a time $t<8\,\mathrm{s}$ and at a displacement $x<16\,\mathrm{m}$. \marks{3} \\
  \item Sketch the velocity--time graph for $P$ from $t=0$ until it first comes to rest. Your sketch should indicate the concavity of the graph, but no numerical values are required. \marks{1} \\
  \item Determine the maximum displacement of $P$ from the origin. \marks{2} \\
\end{enumerate}

\newpage

% ---- Question 5 ----
\questionitem
% [Source: _QBT__Motion_Under_a_Variable_Force, Q5]

A particle $B$ of mass $2\,\mathrm{kg}$ moves along the positive $x$-axis. When $B$ passes through the origin, its speed is $u\,\mathrm{m\,s^{-1}}$. \\

The only horizontal force acting on $B$ is a resistance whose magnitude is directly proportional to the cube of the speed of $B$. When the speed of $B$ is $3\,\mathrm{m\,s^{-1}}$, the magnitude of the resistance is $54\,\mathrm{N}$. \\

\begin{enumerate}[label=\textbf{(\alph*)}]
  \item Show that, while $B$ is moving in the positive $x$-direction, its motion satisfies
  \[
  \frac{1}{v^2}\frac{\mathrm{d}v}{\mathrm{d}x}=-1
  \]
  \marks[-30pt]{3}
  \item
  \begin{enumerate}[label=(\roman*)]
    \item Solve the differential equation in part (a) to obtain $v$ in terms of $x$ and $u$. \marks{4} \\
    \item By considering the behaviour of $v$ as $x\to\infty$, state one respect in which this model of the motion is unrealistic. \marks{1} \\
  \end{enumerate}
\end{enumerate}
The particle subsequently reaches a point $A$ with coordinate $X$, where its speed is $V\,\mathrm{m\,s^{-1}}$. The work done by the resistance as $B$ moves from the origin to $A$ is $W\,\mathrm{J}$. \\
\begin{enumerate}[label=\textbf{(\alph*)}, resume]
  \item
  \begin{enumerate}[label=(\roman*)]
    \item Using $W=\int F\,\mathrm{d}x$, where $F$ is the signed force in the $x$-direction, find an expression for $W$ in terms of $X$ and $u$. \marks{3} \\
    \item Explain the physical significance of the sign of $W$. \marks{1} \\
    \item Rewrite your answer to part (c)(i) in terms of $V$ and $u$, and hence show how $W$ is related to the change in the kinetic energy of $B$. \marks{2} \\
  \end{enumerate}
\end{enumerate}

\newpage

% ---- Question 6 ----
\questionitem
% [Source: _QBT__Motion_Under_a_Variable_Force, Q6]

A small sphere $P$ of mass $5\,\mathrm{kg}$ is released into a liquid and moves vertically downwards. Take the downward direction as positive. In a model of the motion, the forces acting on $P$ are: \\

\begin{itemize}
  \item its weight;
  \item a constant upthrust of magnitude $14\,\mathrm{N}$;
  \item a resistance to motion whose magnitude is proportional to the speed of $P$.
\end{itemize}

When the speed of $P$ is $4\,\mathrm{m\,s^{-1}}$, the resistance has magnitude $10\,\mathrm{N}$. At time $t=0$, the sphere is at a point $O$ and is moving downwards with speed $u\,\mathrm{m\,s^{-1}}$, where $u>0$. At time $t$ seconds, its downward velocity is $v\,\mathrm{m\,s^{-1}}$. Take $g=9.8\,\mathrm{m\,s^{-2}}$. \\

\begin{enumerate}[label=\textbf{(\alph*)}]
  \item Show that, according to the model,
  \[
  \dfrac{\mathrm{d}v}{\mathrm{d}t}=7-\dfrac{1}{2}v
  \]
  \marks[-30pt]{2}
  \item
  \begin{enumerate}[label=(\roman*)]
    \item Solve the differential equation in part (a) to show that
    \[
    v=14+(u-14)e^{-t/2}
    \]
    \marks[-20pt]{5}
    \item Describe briefly how the speed of $P$ varies over time in each of the following cases:
    \begin{itemize}
      \item $0<u<14$;
      \item $u=14$;
      \item $u>14$.
    \end{itemize}
    In each case, state the limiting speed of $P$ as $t\to\infty$. \marks{3} \\
  \end{enumerate}
  \item Given that $u=9$, find the exact downward displacement of $P$ from $O$ when $t=6$. \marks{4} \\
\end{enumerate}

\newpage

% ---- Question 7 ----
\questionitem
% [Source: _QBT__Motion_Under_a_Variable_Force, Q7]

A vehicle of mass $900\,\mathrm{kg}$ moves along a straight horizontal road. Its engine provides a constant driving force of $576\,\mathrm{N}$. When the vehicle has speed $v\,\mathrm{m\,s^{-1}}$, the resistance to its motion has magnitude $(v^2+14v)\,\mathrm{N}$. \\

The time taken for the speed of the vehicle to increase from $3\,\mathrm{m\,s^{-1}}$ to $13\,\mathrm{m\,s^{-1}}$ is $T$ seconds. \\

\begin{enumerate}[label=\textbf{(\alph*)}]
  \item Show that
  \[
  T=18\ln\left(\frac{27}{7}\right)
  \]
  \marks[-30pt]{7}
\end{enumerate}
When $t=0$, the speed of the vehicle is $8\,\mathrm{m\,s^{-1}}$. \\
\begin{enumerate}[label=\textbf{(\alph*)}, resume]
  \item Show that
  \[
  \frac{v+32}{18-v}=pe^{qt}
  \]
  where the rational numbers $p$ and $q$ are to be found. \marks{2} \\
  \item Hence explain why $v<18$ for all $t\geq 0$. \marks{1} \\
\end{enumerate}

\newpage

% ---- Question 8 ----
\questionitem
% [Source: _QBT__Motion_Under_a_Variable_Force, Q8]

A particle $P$ of mass $2\,\mathrm{kg}$ moves along the positive $x$-axis in the direction of increasing $x$. At time $t$ seconds, its displacement from the origin is $x$ metres and its speed is $v\,\mathrm{m\,s^{-1}}$. \\

While $P$ is moving, the resultant force acting on it is directed towards the origin and has magnitude $kv(v+2)\,\mathrm{N}$, where $k$ is a positive constant and $v$ denotes the numerical value of the speed in metres per second. \\

When $x=1$, the speed of $P$ is $7\,\mathrm{m\,s^{-1}}$, and when $x=3$, its speed is $1\,\mathrm{m\,s^{-1}}$. \\

\begin{enumerate}[label=\textbf{(\alph*)}]
  \item Find $k$ and show that
  \[
  v=9\cdot 3^{(1-x)/2}-2
  \]
  \marks[-20pt]{6}
\end{enumerate}
The time taken for the speed of $P$ to decrease from $7\,\mathrm{m\,s^{-1}}$ to $1\,\mathrm{m\,s^{-1}}$ is $T$ seconds. \\
\begin{enumerate}[label=\textbf{(\alph*)}, resume]
  \item Show that
  \[
  T=\frac{\ln(7/3)}{\ln 3}
  \]
  \marks[-30pt]{4}
\end{enumerate}

\newpage

% ---- Question 9 ----
\questionitem
% [Source: _QBT__Motion_Under_a_Variable_Force, Q9]

A particle $P$ of mass $4\,\mathrm{kg}$ is released from rest and moves vertically downwards through a fluid. While $P$ is moving with speed $v\,\mathrm{m\,s^{-1}}$, it experiences a resistance of magnitude $6v\,\mathrm{N}$ acting vertically upwards. After $P$ has descended a distance $x\,\mathrm{m}$, its speed is $v\,\mathrm{m\,s^{-1}}$. Take $g=9.8\,\mathrm{m\,s^{-2}}$. \\

\begin{enumerate}[label=\textbf{(\alph*)}]
  \item Show that
  \[
  \left(\frac{10v}{98-15v}\right)\frac{\mathrm{d}v}{\mathrm{d}x}=1
  \]
  \marks[-30pt]{2}
  \item Verify that $x$ and $v$ satisfy
  \[
  45x=196\ln\left(\frac{98}{98-15v}\right)-30v
  \]
  \marks[-30pt]{4}
  \item Determine the work done against the resistance from the instant $P$ is released until its acceleration has magnitude $2.45\,\mathrm{m\,s^{-2}}$. Give your answer to $3$ significant figures. \marks{5} \\
\end{enumerate}

\newpage

% ---- Question 10 ----
\questionitem
% [Source: _QBT__Motion_Under_a_Variable_Force, Q10]

A particle $P$, of mass $m$, moves on a rough straight track inclined at an angle $\alpha$ to the horizontal. A fixed point $O$ lies on the track, and $x$ denotes the distance of $P$ from $O$, measured up the track. \\

In addition to its weight, $P$ is acted on by a force directed down the track and having magnitude
\[
\frac{kmg}{x+a}
\] \\
where $a$ and $k$ are positive constants with dimensions of length. The coefficient of friction between $P$ and the track is $\mu$. \\

The particle is projected up the track and first comes to rest at the point where $x=2a$. \\

\begin{enumerate}[label=\textbf{(\alph*)}]
  \item Show that, while $P$ is moving up the track with speed $v$,
  \[
  \frac{\mathrm{d}}{\mathrm{d}x}(v^2)+\frac{2gk}{x+a}+2g(\sin\alpha+\mu\cos\alpha)=0
  \]
  \marks[-30pt]{3}
  \item Verify that
  \[
  v^2=2gk\ln\left(\frac{3a}{x+a}\right)+2g(\sin\alpha+\mu\cos\alpha)(2a-x)
  \]
  \marks[-30pt]{3}
  \item Find, in terms of $a$, $k$ and $\alpha$, the range of values of $\mu$ for which $P$ remains at rest when $x=2a$. \marks{2} \\
\end{enumerate}

\newpage

% ---- Question 11 ----
\questionitem
% [Source: _QBT__Motion_Under_a_Variable_Force, Q11]

A package of mass $m\,\mathrm{kg}$ is released from rest and falls vertically. When its downward speed is $v\,\mathrm{m\,s^{-1}}$, it experiences a resistive force of magnitude $mkv^2\,\mathrm{N}$, where $k$ is a positive constant. Take the downward direction as positive. \\

\begin{enumerate}[label=\textbf{(\alph*)}]
  \item Show that, at time $t\,\mathrm{s}$ after release,
  \[
  v=\sqrt{\frac{g}{k}}\left(\frac{e^{2\sqrt{gk}\,t}-1}{e^{2\sqrt{gk}\,t}+1}\right)
  \]
  \marks[-40pt]{4}
\end{enumerate}
It is given that $k=0.0125\,\mathrm{m^{-1}}$. \\
\begin{enumerate}[label=\textbf{(\alph*)}, resume]
  \item Find the distance fallen by the package during the first $2$ seconds after its release. \marks{4} \\
\end{enumerate}

\newpage

% ---- Question 12 ----
\questionitem
% [Source: _QBT__Motion_Under_a_Variable_Force, Q12]

A particle $P$ of mass $16\,\mathrm{kg}$ moves in a straight line on a smooth horizontal plane. \\ 

At time $t=0$, the particle passes through a fixed point $O$ with velocity $\dfrac{5\sqrt{3}}{4}\,\mathrm{m\,s^{-1}}$. \\

From this instant, a horizontal force opposing the motion acts on $P$. The magnitude of the force is proportional to $\sqrt{25-4v^2}$, where $v\,\mathrm{m\,s^{-1}}$ is the velocity of $P$.  \\

When $v=2$, the magnitude of the force is $6\,\mathrm{N}$. The force ceases to act when the velocity has decreased to $\dfrac{5}{4}\,\mathrm{m\,s^{-1}}$. \\

\begin{enumerate}[label=\textbf{(\alph*)}]
  \item Show that, while the force is acting, the motion of $P$ is modelled by the differential equation
  \[
  \frac{1}{\sqrt{25-4v^2}}\frac{\mathrm{d}v}{\mathrm{d}t}=-\frac{1}{8}
  \]
  \marks[-30pt]{3}
  \item Hence find $v$ in terms of $t$, showing that your expression is valid for
  \[
  0\leq t\leq \frac{2\pi}{3}
  \]
  \marks[-20pt]{5}
  \item Determine the total distance travelled by $P$ during the interval $0\leq t\leq 8$. Give your answer in an exact form. \marks{4} \\
\end{enumerate}

\newpage

% ---- Question 13 ----
\questionitem
% [Source: _QBT__Motion_Under_a_Variable_Force, Q13]

A small instrument package of mass $m\,\mathrm{kg}$ enters a deep liquid moving vertically downwards at $7\,\mathrm{m\,s^{-1}}$. While the package is moving downwards, it experiences a resistive force of magnitude $\dfrac{m}{20}v^2\,\mathrm{N}$, where $v\,\mathrm{m\,s^{-1}}$ is its speed. \\

At time $t\,\mathrm{s}$ after entering the liquid, the package has descended a distance $x\,\mathrm{m}$. \\

\begin{enumerate}[label=\textbf{(\alph*)}]
  \item Find an expression for $v$ in terms of $t$. \marks{6} \\
  \item Find an expression for $x$ in terms of $t$. \marks{3} \\
  \item Find the exact distance descended by the package when its speed first reaches $\dfrac{21}{2}\,\mathrm{m\,s^{-1}}$. \marks{2} \\
\end{enumerate}

\newpage

% ---- Question 14 ----
\questionitem
% [Source: _QBT__Motion_Under_a_Variable_Force, Q14]

A vehicle of mass $m\,\mathrm{kg}$ moves along a horizontal straight road. Its engine works at a constant rate of $2kmg\,\mathrm{W}$, where $k$ is a positive constant with units of speed. When the vehicle has speed $v\,\mathrm{m\,s^{-1}}$, it experiences a resistance of magnitude $\dfrac{mgv}{k}\,\mathrm{N}$. \\

Initially, the vehicle is at rest at $x=0$, where $x\,\mathrm{m}$ is the distance travelled. \\

\begin{enumerate}[label=\textbf{(\alph*)}]
  \item Show that
  \[
  v^2\dfrac{\mathrm{d}v}{\mathrm{d}x}=\dfrac{g}{k}\left(2k^2-v^2\right)
  \]
  \marks[-30pt]{3}
  \item Verify that
  \[
  gx=-kv+\dfrac{k^2}{\sqrt{2}}\ln\left(\dfrac{\sqrt{2}k+v}{\sqrt{2}k-v}\right)
  \]
  \marks[-40pt]{5}
  \item By using the work--energy principle, show that the time taken for the vehicle to accelerate from rest to a speed $V$, where $0<V<\sqrt{2}k$, is
  \[
  \dfrac{k}{2g}\ln\left(\dfrac{2k^2}{2k^2-V^2}\right)
  \]
  \marks[-40pt]{4}
\end{enumerate}

\end{enumerate}
\end{document}
